%% ma: L12-B14-0037
%% lop: 12
%% chuong: 5
%% bai: 14
%% dang: 12.14.19
%% loai: TN4
%% muc: VD
%% dap_an: "A"
%% nguon: "CĐ2.5.pdf giáo viên gửi 27/10/2026 (M07, Lê Minh Khôi – Nguyễn Anh Khôi, Chương II Chủ đề 5 Tích có hướng của hai vectơ), Câu 21"
%% kien_thuc: "Bài 14 (tích có hướng, phương trình mặt phẳng, khoảng cách từ điểm đến mặt phẳng), lớp 11 (thể tích khối chóp, khối hộp)"
%% trang_thai: cho-duyet
%% ly_do: "Dạng mới 12.14.19 đề xuất, chờ giáo viên duyệt"
\begin{ex}
Trong không gian $Oxyz$, cho tứ diện $ABCD$ có $A(2;-1;1)$, $B(3;0;-1)$, $C(2;-1;3)$ và điểm $D$ di động trên trục $Oy$. Biết rằng tứ diện $ABCD$ có thể tích bằng $5$. Tính tổng tung độ của các điểm $D$ thỏa mãn.
\choice{\True $-6$}{$2$}{$7$}{$-4$}
\loigiai{$D(0;y;0)$, $\vec n=\left[\overrightarrow{AB},\overrightarrow{AC}\right]=(2;-2;0)$, $\vec n\cdot\overrightarrow{AD}=-2y-6$.

$V=\dfrac{|2y+6|}6=5\Leftrightarrow y=12$ hoặc $y=-18$. Tổng bằng $-6$.}
\end{ex}
